% Options for packages loaded elsewhere
\PassOptionsToPackage{unicode}{hyperref}
\PassOptionsToPackage{hyphens}{url}
%
\documentclass[
]{article}
\usepackage{amsmath,amssymb}
\usepackage{iftex}
\ifPDFTeX
  \usepackage[T1]{fontenc}
  \usepackage[utf8]{inputenc}
  \usepackage{textcomp} % provide euro and other symbols
\else % if luatex or xetex
  \usepackage{unicode-math} % this also loads fontspec
  \defaultfontfeatures{Scale=MatchLowercase}
  \defaultfontfeatures[\rmfamily]{Ligatures=TeX,Scale=1}
\fi
\usepackage{lmodern}
\ifPDFTeX\else
  % xetex/luatex font selection
\fi
% Use upquote if available, for straight quotes in verbatim environments
\IfFileExists{upquote.sty}{\usepackage{upquote}}{}
\IfFileExists{microtype.sty}{% use microtype if available
  \usepackage[]{microtype}
  \UseMicrotypeSet[protrusion]{basicmath} % disable protrusion for tt fonts
}{}
\makeatletter
\@ifundefined{KOMAClassName}{% if non-KOMA class
  \IfFileExists{parskip.sty}{%
    \usepackage{parskip}
  }{% else
    \setlength{\parindent}{0pt}
    \setlength{\parskip}{6pt plus 2pt minus 1pt}}
}{% if KOMA class
  \KOMAoptions{parskip=half}}
\makeatother
\usepackage{xcolor}
\usepackage[margin=1in]{geometry}
\usepackage{longtable,booktabs,array}
\usepackage{calc} % for calculating minipage widths
% Correct order of tables after \paragraph or \subparagraph
\usepackage{etoolbox}
\makeatletter
\patchcmd\longtable{\par}{\if@noskipsec\mbox{}\fi\par}{}{}
\makeatother
% Allow footnotes in longtable head/foot
\IfFileExists{footnotehyper.sty}{\usepackage{footnotehyper}}{\usepackage{footnote}}
\makesavenoteenv{longtable}
\setlength{\emergencystretch}{3em} % prevent overfull lines
\providecommand{\tightlist}{%
  \setlength{\itemsep}{0pt}\setlength{\parskip}{0pt}}
\setcounter{secnumdepth}{-\maxdimen} % remove section numbering
\ifLuaTeX
  \usepackage{selnolig}  % disable illegal ligatures
\fi
\IfFileExists{bookmark.sty}{\usepackage{bookmark}}{\usepackage{hyperref}}
\IfFileExists{xurl.sty}{\usepackage{xurl}}{} % add URL line breaks if available
\urlstyle{same}
\hypersetup{
  hidelinks,
  pdfcreator={LaTeX via pandoc}}

\author{}
\date{}

\begin{document}

\hypertarget{robocam-stentor-photobehavior-assay-protocol}{%
\section{RoboCam Stentor Photobehavior Assay
Protocol}\label{robocam-stentor-photobehavior-assay-protocol}}

Adapted from a design document our PI provided (v1.0, July 2026; kept in
a local reference library, not committed here for copyright reasons).
This is our own restructured summary of that protocol, cross-referenced
against
\href{WAVELENGTH_SELECTION_GUIDE.md}{\texttt{WAVELENGTH\_SELECTION\_GUIDE.md}},
which it independently corroborates: primary stimulus 610/620 nm,
secondary 480/540 nm, and \textbf{680 nm explicitly as a
comparison/control wavelength, not the main stimulus} --- the same
conclusion our own literature synthesis reached.

\hypertarget{the-three-behaviors-one-experiment-each}{%
\subsection{The three behaviors, one experiment
each}\label{the-three-behaviors-one-experiment-each}}

\begin{longtable}[]{@{}
  >{\raggedright\arraybackslash}p{(\columnwidth - 8\tabcolsep) * \real{0.2000}}
  >{\raggedright\arraybackslash}p{(\columnwidth - 8\tabcolsep) * \real{0.2000}}
  >{\raggedright\arraybackslash}p{(\columnwidth - 8\tabcolsep) * \real{0.2000}}
  >{\raggedright\arraybackslash}p{(\columnwidth - 8\tabcolsep) * \real{0.2000}}
  >{\raggedright\arraybackslash}p{(\columnwidth - 8\tabcolsep) * \real{0.2000}}@{}}
\toprule\noalign{}
\begin{minipage}[b]{\linewidth}\raggedright
Behavior
\end{minipage} & \begin{minipage}[b]{\linewidth}\raggedright
Core question
\end{minipage} & \begin{minipage}[b]{\linewidth}\raggedright
Stimulus geometry
\end{minipage} & \begin{minipage}[b]{\linewidth}\raggedright
Primary metric
\end{minipage} & \begin{minipage}[b]{\linewidth}\raggedright
Best RoboCam mode
\end{minipage} \\
\midrule\noalign{}
\endhead
\bottomrule\noalign{}
\endlastfoot
Step-up photophobic response & Does a sudden light increase trigger
stop/reversal/turn? & Whole-field temporal pulse, or dark-to-light edge
& Response probability, latency, reversal duration, turn angle &
High-frame-rate single-well event assay \\
Phototaxis & Does \emph{Stentor} move directionally relative to a light
field? & Half-field, side illumination, or smooth gradient along the
light axis & Occupancy index and velocity projection & Flat chamber,
stable spatial light field \\
Photokinesis & Does light alter swimming speed independent of direction?
& Uniform full-field illumination at a defined intensity & Speed ratio,
light vs.~baseline & 24-well screen with replicated wells \\
\end{longtable}

\textbf{Recommended order: photophobic response first, then phototaxis,
then photokinesis.} Phototaxis results are ambiguous on their own ---
dark-side accumulation in a half-field assay can be produced by
edge-triggered avoidance (photophobic response) rather than true
directional steering, so it needs the photophobic assay validated first
as a baseline to interpret against.

\hypertarget{design-principles}{%
\subsection{Design principles}\label{design-principles}}

\begin{itemize}
\tightlist
\item
  \textbf{Separate stimulus geometry from behavior class.} Temporal
  steps for photophobia, gradients/half-fields for taxis, uniform light
  for kinesis. Conflating these risks mislabeling a photophobic edge
  artifact as ``phototaxis.''
\item
  \textbf{Synchronize video and LED timing.} Record LED-on timestamps in
  the same data stream as the video frame index (or add a visible
  timing-indicator LED). Without this, latency and event scoring can't
  be trusted.
\item
  \textbf{Use shallow, optically flat chambers.} Flatten the meniscus,
  constrain depth, keep \emph{Stentor} in a single focal layer ---
  losing focus through depth is usually a bigger problem than
  insufficient magnification, since \emph{Stentor} is already large
  (approaching/exceeding 1 mm extended).
\item
  \textbf{Calibrate irradiance at the sample plane.} Report wavelength,
  irradiance, pulse duration, duty cycle, and chamber geometry --- PWM
  percentage or supply voltage alone isn't comparable across builds.
\item
  \textbf{Use non-actinic imaging light.} Dim IR (850/940 nm) or very
  dim red/far-red for tracking, verified with a tracking-light-only
  control to confirm it doesn't itself elicit a response.
\item
  \textbf{Block and randomize.} Randomize well order, wavelength order,
  and intensity sequence within culture-day blocks, so
  time/evaporation/fatigue/culture condition doesn't masquerade as a
  light effect.
\end{itemize}

\hypertarget{stimulus-wavelengths}{%
\subsection{Stimulus wavelengths}\label{stimulus-wavelengths}}

\begin{longtable}[]{@{}
  >{\raggedright\arraybackslash}p{(\columnwidth - 4\tabcolsep) * \real{0.3333}}
  >{\raggedright\arraybackslash}p{(\columnwidth - 4\tabcolsep) * \real{0.3333}}
  >{\raggedright\arraybackslash}p{(\columnwidth - 4\tabcolsep) * \real{0.3333}}@{}}
\toprule\noalign{}
\begin{minipage}[b]{\linewidth}\raggedright
Channel
\end{minipage} & \begin{minipage}[b]{\linewidth}\raggedright
Purpose
\end{minipage} & \begin{minipage}[b]{\linewidth}\raggedright
Notes
\end{minipage} \\
\midrule\noalign{}
\endhead
\bottomrule\noalign{}
\endlastfoot
610 or 620 nm & Primary stimulus & Most important channel for
photophobic and phototaxis assays \\
540 or 480 nm & Secondary action-spectrum point & Useful after the main
assay works; don't start here \\
660/680 nm & Long-red comparison & Expected to be less effective than
610/620 nm; used as a control \\
850/940 nm & Imaging only & Verify with a tracking-light-only control \\
\end{longtable}

Use a constant-current LED driver rather than driving high-power LEDs
directly from a microcontroller; have RoboCam control the driver enable
pin or a logic-level MOSFET gate, and log the actual command waveform
(ideally alongside a photodiode signal from the sample plane during
validation).

\hypertarget{chamber-design}{%
\subsection{Chamber design}\label{chamber-design}}

\begin{itemize}
\tightlist
\item
  \textbf{First assay:} a shallow chamber, not an open 24-well meniscus
  --- glass-bottom or optically clear well with a removable
  spacer/gasket and a top cover film/coverslip. Target depth 0.5-1.5 mm
  (deep enough to avoid compressing cells, shallow enough for focus and
  intensity uniformity); planar area 8-15 mm for single-well tracking.
\item
  \textbf{24-well plates} are appropriate for photokinesis screens and
  replicated intensity/wavelength tests, but raw deep-well distributions
  shouldn't be interpreted as taxis unless the light field and depth are
  characterized --- curved meniscus and wall effects can dominate the
  distribution.
\item
  For one-well-at-a-time scanning, randomize well order so plate
  position and elapsed time aren't confounded with condition.
\end{itemize}

\hypertarget{assay-1-step-up-photophobic-response-start-here}{%
\subsection{Assay 1 --- Step-up photophobic response (start
here)}\label{assay-1-step-up-photophobic-response-start-here}}

Recommended as the first RoboCam validation assay: gives a time-locked
event with clear behavioral labels.

\begin{itemize}
\tightlist
\item
  \textbf{Pulse durations:} 0.2, 0.5, 1.0, 2.0 s. \textbf{Intensity:}
  log-spaced pilot series at the sample plane, to fit dose-response
  curves and estimate an EC50-like threshold. \textbf{Inter-trial
  interval:} at least 5-10 s (longer if repeated stimulation shows
  adaptation). \textbf{Frame rate:} 60 fps minimum, 120 fps preferred
  for latency/reversal dynamics.
\item
  \textbf{Event definitions:} baseline speed = median speed -2.0 to -0.2
  s before LED onset; response window = 0.05-1.5 s after onset; stop
  event = speed drops below 20-30\% of baseline for \textgreater= 2-3
  consecutive frames; reversal event = negative displacement relative to
  pre-stimulus heading for \textgreater= 100-150 ms; latency = time from
  LED onset to first stop/reversal frame; turn angle = angle between
  mean heading before stimulus and mean heading after recovery.
\item
  \textbf{Output metrics:} response probability, median latency/IQR by
  wavelength and irradiance, reversal duration and backward distance,
  turn-angle distribution, dose-response curve (probability vs.~log
  irradiance), adaptation curve across repeated pulses.
\end{itemize}

\hypertarget{assay-2-phototaxis-directional-bias-only-after-assay-1-works}{%
\subsection{Assay 2 --- Phototaxis / directional bias (only after Assay
1
works)}\label{assay-2-phototaxis-directional-bias-only-after-assay-1-works}}

\begin{itemize}
\tightlist
\item
  \textbf{Half-field chamber:} mask/project half the chamber dark, half
  bright at 610/620 nm. Readout: dark-side occupancy over time and
  edge-crossing response --- but this conflates true steering with
  photophobic boundary avoidance.
\item
  \textbf{Smooth gradient:} diffuser + neutral-density gradient, wedge
  mask, or side illumination across the chamber. Readout: velocity
  projection along the gradient axis --- better isolates steering/taxis,
  but needs a calibrated light field.
\item
  \textbf{Metrics:} occupancy index = (N\_dark - N\_light)/N\_total,
  reported over time, not just at the end; directional bias = mean
  velocity projected away from the light source, divided by speed (near
  +1 = moving away, near 0 = no bias); edge-crossing response
  probability; dwell time per zone, normalized by zone area.
\item
  \textbf{Caution:} don't call a result ``true phototaxis'' from
  dark-side accumulation alone --- pair it with velocity-vector analysis
  and the photokinesis control. Avoid bubbles, curved meniscus, and wall
  shadows, which generate artificial gradients and trapping zones.
\end{itemize}

\hypertarget{assay-3-photokinesis-speed-modulation}{%
\subsection{Assay 3 --- Photokinesis / speed
modulation}\label{assay-3-photokinesis-speed-modulation}}

The best fit for 24-well RoboCam scans, since uniform full-field
illumination is easier to reproduce across wells than a gradient.

\begin{itemize}
\tightlist
\item
  \textbf{Windows:} baseline 30-60 s before light; acute onset 0-2 s
  after light-on (inspect, but don't use as the main metric ---
  photophobic events can contaminate speed here); sustained light 10-60
  s after light-on (primary photokinesis estimate); recovery 30-60 s
  after light-off (reversibility/fatigue).
\item
  \textbf{Photokinesis index} = median speed during sustained light /
  median baseline speed. Use a paired design --- each well/cell as its
  own baseline control. Exclude tracks that hit walls, attach to
  substrate, or leave focus for more than a defined fraction of the
  window.
\item
  \textbf{First screen:} 610/620 nm at low/medium/high irradiance, plus
  a no-light control and a 680 nm comparison.
\end{itemize}

\hypertarget{controls-and-confound-checks}{%
\subsection{Controls and confound
checks}\label{controls-and-confound-checks}}

\begin{longtable}[]{@{}
  >{\raggedright\arraybackslash}p{(\columnwidth - 4\tabcolsep) * \real{0.3333}}
  >{\raggedright\arraybackslash}p{(\columnwidth - 4\tabcolsep) * \real{0.3333}}
  >{\raggedright\arraybackslash}p{(\columnwidth - 4\tabcolsep) * \real{0.3333}}@{}}
\toprule\noalign{}
\begin{minipage}[b]{\linewidth}\raggedright
Control
\end{minipage} & \begin{minipage}[b]{\linewidth}\raggedright
Purpose
\end{minipage} & \begin{minipage}[b]{\linewidth}\raggedright
Pass criterion
\end{minipage} \\
\midrule\noalign{}
\endhead
\bottomrule\noalign{}
\endlastfoot
No-actinic sham & Camera/vibration/software-trigger artifacts & No
increase in stop/reversal events at sham time \\
Tracking light only & Confirms IR/far-red imaging is behaviorally silent
& Baseline speed/event frequency match dark control \\
680 nm comparison & Separates wavelength-specific sensitivity from
generic brightness/heating & Lower or shifted response vs.~610/620 nm at
matched irradiance \\
Heat control & Tests whether LED warming changes swimming speed &
Temperature rise below a predefined limit; speed not explained by heat
alone \\
Spatial mask blank & Tests whether masks create flow/shadow/focusing
artifacts & No preferential accumulation without actinic contrast \\
Fixed particle video & Detects stage drift and optical distortion & No
apparent movement in fixed particles during LED switching \\
\end{longtable}

\hypertarget{recommended-workflow}{%
\subsection{Recommended workflow}\label{recommended-workflow}}

\textbf{Phase 0 --- build and calibration:} measure field uniformity
(photodiode/camera flat-field per mask and LED channel); verify LED
timing (film an indicator LED or log a photodiode signal alongside
video); measure sample-plane temperature during the longest planned
exposure; run tracking-light-only controls; record a calibration video
with beads/fixed particles for pixel-to-micron scale.

\textbf{Phase 1 --- minimum viable experiment:} load 5-20 healthy
swimming cells into a shallow chamber (avoid dense cultures, disturbed
cells, debris); 5-10 min acclimation/dark adaptation; 10 s baseline
under non-actinic light; a 610/620 nm pulse series (0.5 s and 1.0 s at
low/medium/high irradiance, 5-10 s between trials); manually score the
first 20-50 events before comparing to automated scoring; repeat with
sham pulses and a 680 nm comparison at matched irradiance.

\textbf{Phase 2 --- publication-grade assay:} \textgreater= 3
independent culture days/batches as biological replicates; multiple
wells/chambers per condition per day as technical replicates; report
analyzable cells/tracks, but treat culture day/well as the primary
replicate level; randomize intensity and well order (avoid monotonic
low-to-high sequences unless specifically studying adaptation); blind
manual event scoring to wavelength/intensity where practical; predefine
QC exclusions (stuck cells, collisions, track loss, fission-stage cells,
debris).

\hypertarget{analysis-pipeline}{%
\subsection{Analysis pipeline}\label{analysis-pipeline}}

Preprocessing (flat-field correction, background subtraction, timestamp
alignment) -\textgreater{} segmentation (contour/threshold tuned for
large \emph{Stentor} bodies, ignoring sub-threshold particles)
-\textgreater{} tracking (nearest-neighbor/Kalman linking with
gap-closing) -\textgreater{} event detection (speed/heading/displacement
relative to the pre-stimulus vector) -\textgreater{} metrics (response
probability, latency, reversal duration, speed ratio, occupancy index,
directional bias) -\textgreater{} archiving (raw video, metadata
JSON/CSV, processed tracks, analysis code version --- a reproducible
experiment folder per run).

\hypertarget{minimum-metadata-to-record-per-run}{%
\subsection{Minimum metadata to record per
run}\label{minimum-metadata-to-record-per-run}}

Experiment (id, date, operator, RoboCam build/version, analysis
version); biology (species, culture source/age, feeding status, medium,
temperature); chamber (plate type, depth, gasket material, volume, cover
method, well ID); optics (camera model, lens, magnification, frame rate,
resolution, pixel size, imaging light wavelength/intensity); stimulus
(LED wavelength, bandwidth if known, irradiance at sample plane, pulse
duration, duty cycle, mask/gradient geometry); trial (baseline duration,
stimulus onset frame, recovery duration, inter-trial interval,
randomized order index); QC (cells loaded, cells analyzable, exclusions,
notes on bubbles/debris/focus).

\hypertarget{statistics}{%
\subsection{Statistics}\label{statistics}}

\begin{itemize}
\tightlist
\item
  Photophobic response probability: logistic regression / generalized
  mixed model, wavelength/intensity as fixed effects, culture day/well
  as random/blocking factors.
\item
  Latency and reversal duration: medians or robust models (distributions
  are often skewed and censored by frame rate).
\item
  Photokinesis: paired speed ratios by cell or well, aggregated by
  well/culture day --- don't treat every frame as an independent
  replicate.
\item
  Phototaxis: report both occupancy and directional movement ---
  occupancy alone is insufficient, since cells can accumulate in dark
  zones through photophobic boundary rejection rather than true
  steering.
\item
  Report failed tracking and exclusions transparently; these assays are
  highly sensitive to chamber quality and culture condition.
\end{itemize}

\hypertarget{suggested-24-well-plate-layout-for-the-initial-screen}{%
\subsection{Suggested 24-well plate layout for the initial
screen}\label{suggested-24-well-plate-layout-for-the-initial-screen}}

\begin{longtable}[]{@{}
  >{\raggedright\arraybackslash}p{(\columnwidth - 4\tabcolsep) * \real{0.3333}}
  >{\raggedright\arraybackslash}p{(\columnwidth - 4\tabcolsep) * \real{0.3333}}
  >{\raggedright\arraybackslash}p{(\columnwidth - 4\tabcolsep) * \real{0.3333}}@{}}
\toprule\noalign{}
\begin{minipage}[b]{\linewidth}\raggedright
Wells
\end{minipage} & \begin{minipage}[b]{\linewidth}\raggedright
Condition
\end{minipage} & \begin{minipage}[b]{\linewidth}\raggedright
Purpose
\end{minipage} \\
\midrule\noalign{}
\endhead
\bottomrule\noalign{}
\endlastfoot
A1-A3 & No-actinic sham, imaging only & Baseline / tracking-light
control \\
A4-A6 & 610/620 nm low irradiance & Lower end of dose-response \\
B1-B3 & 610/620 nm medium irradiance & Expected working condition \\
B4-B6 & 610/620 nm high irradiance & Saturation/adaptation check \\
C1-C3 & 680 nm matched irradiance & Long-red wavelength comparison \\
C4-C6 & Uniform 610/620 nm sustained light & Photokinesis window \\
D1-D3 & Half-field 610/620 nm & Phototaxis/edge-avoidance pilot \\
D4-D6 & Recovery/dark after exposure & Carryover and fatigue check \\
\end{longtable}

Don't run the plate in A1-to-D6 order for one-well-at-a-time scanning
--- generate a randomized acquisition order instead.

\end{document}
